\documentclass[conference]{IEEEtran}

% ── Packages ────────────────────────────────────────────────────────────────
\usepackage[T1]{fontenc}
\usepackage[utf8]{inputenc}
\usepackage{amsmath,amssymb}
\usepackage{graphicx}
\usepackage{booktabs}
\usepackage{array}
\usepackage{multirow}
\usepackage{cite}
\usepackage{url}
\usepackage{xcolor}
\usepackage{hyperref}
\hypersetup{colorlinks=false, hidelinks}

% ── Metadata ─────────────────────────────────────────────────────────────────
\title{DisasterAI: A Predictive Multi-Agent Reinforcement Learning\\
Framework for Real-Time Urban Flood Emergency Response}

\author{
  \IEEEauthorblockN{Viraj H.~C.\textsuperscript{1},\quad Abhinav Tripathi\textsuperscript{1}}
  \IEEEauthorblockA{
    \textit{Supervisor: Mohandas R.}\\
    \textit{Department of Computer Science \& Engineering (AIML)}\\
    SRMIST, Kattankulathur, Chennai, India\\
    April 2026
  }
}

% ════════════════════════════════════════════════════════════════════════════
\begin{document}
\maketitle

% ── Abstract ─────────────────────────────────────────────────────────────────
\begin{abstract}
Traditional reactive dispatch systems break down during urban flooding because
roads become impassable faster than field units can reroute.
This paper presents \textbf{DisasterAI}, a simulation framework that combines a
predictive Multi-Agent Reinforcement Learning (MARL) architecture with a D8
priority-queue flood engine, risk-weighted Hungarian dispatch, and
predictive flood-aware A* pathfinding.
The simulation targets a monsoon flood in Mumbai's Bandra-Kurla~Complex /
Mithi River basin using NASA SRTM elevation data, OpenStreetMap road networks,
live Open-Meteo discharge API feeds, and GDACS severity alerts.
A composite victim risk score---current depth, predicted future depth,
health-decay pressure, and population vulnerability---drives all dispatch
and routing decisions.
Evaluated against five baseline strategies over twenty independent runs,
Hungarian dispatch achieves a mean simulation score of $763.0$ versus $202.9$
for Greedy and $-28.2$ for Priority Queue baselines, with mean response time
19.2~steps versus 30.2 for Greedy, a 36.4\% reduction.
Predictive A* routing keeps approximately 40\% of fleet units out of flooded
road segments compared to static routing.
Ablation studies show a lookahead of $N=2$ steps maximises simulation score
($877.9 \pm 396.3$).
The QMIX agent learns to spread across the study area; independent
Q-learners cluster instead.
\end{abstract}

\begin{IEEEkeywords}
Multi-Agent Reinforcement Learning, QMIX, flood simulation, emergency dispatch,
A* pathfinding, Hungarian algorithm, MCLP, digital elevation model,
priority-flood, urban disaster response
\end{IEEEkeywords}

% ════════════════════════════════════════════════════════════════════════════
\section{Introduction}
\label{sec:intro}

The 2005 Mumbai floods deposited 944~mm of rain on the Bandra-Kurla Complex
in 24~hours and killed more than 1,000 people.
A large part of the death toll traces back to dispatch failure: ambulances
routed into already-flooded streets, multiple units converging on the same
accessible areas while isolated victims received no coverage, and no mechanism
to anticipate where roads would be passable by the time a unit arrived.

Static routing systems compute paths at departure time, not at arrival time.
In a fast-moving flood that gap can be lethal.
A route that is clear at 09:00 may be under half a metre of water by 09:20,
which is precisely when the unit gets there.
Addressing this requires a system that runs flood physics forward in time and
uses the result when making routing and dispatch decisions.

This paper describes DisasterAI, a Python simulation built on real Mumbai
geodata that attacks this problem with four components working together:
a D8 priority-queue flood engine that runs predictive forward passes,
a composite risk scorer that weights future flood depth at a higher coefficient
than current depth, a Hungarian dispatch module that operates on that risk
score rather than raw distance, and a QMIX multi-agent policy trained to
coordinate regional coverage.

Our contributions are:
\begin{enumerate}
  \item An open simulation environment for urban monsoon flood response using
    real DEM, OSM, and live API data for the Bandra-Kurla~/ Mithi River basin.
  \item A predictive flood engine that runs D8 propagation forward to the
    unit's estimated arrival time, time-locking routing constraints to arrival
    conditions rather than departure conditions.
  \item A composite victim risk formula whose dominant term ($\beta = 0.40$)
    is predicted future flood depth, validated through ablation.
  \item Baseline and ablation evaluations across 20 runs showing Hungarian
    dispatch scores $3.76\times$ higher than the next-best strategy, and that
    $N=2$ lookahead steps maximise cumulative reward.
\end{enumerate}

% ════════════════════════════════════════════════════════════════════════════
\section{Related Work}
\label{sec:related}

\subsection{Flood Simulation and Digital Elevation Models}

Barnes et al.~\cite{barnes2014priority} established the Priority-Flood
algorithm, which fills DEM depressions by processing cells in elevation order
via a priority queue, achieving O$(n \log n)$ complexity.
Our \texttt{HazardPropagation} class implements this directly.
Ma et al.~\cite{ma2025evaluation} benchmark six queue structures within
Priority-Flood and confirm the min-heap is efficient at DEM scale.
A hash-heap variant~\cite{ma2026modified} gains 2--25\% speed, and the
DZFlood dual-queue~\cite{wu2025efficient} reports 19--28\% improvements on
large grids---both showing the paradigm is still being pushed forward.
GraphFlood~\cite{gailleton2024graphflood} independently validates the
graph-theoretic approach, reaching $10\times$ speedups over shallow-water
equation solvers without sacrificing physical accuracy.
Barnes et al.~\cite{barnes2020hierarchies} explain why min-heap flow correctly
models basin pooling before ridgeline overflow, which is exactly the behaviour
we rely on for the Mithi River basin topography.

Yang et al.~\cite{yang2024cnn} and Wijaya and Yang~\cite{wijaya2021combining}
document why cellular automata only achieve physical plausibility with
deep-learning augmentation---overhead the priority-queue approach avoids.

\subsection{Emergency Dispatch and Task Allocation}

Lee and Lee~\cite{lee2020marl} formalise disaster response as a
partially-observable MARL problem; their cooperative agent framework is
the direct precursor to ours.
For task allocation, bipartite matching is well established:
CF-HMRTA~\cite{bachler2022cfhmrta} handles coalition formation for
heterogeneous robot fleets, while the GRL bigraph framework~\cite{prorok2023bigraph}
augments matching with learned incentive functions.
Sivagnanam et al.~\cite{sivagnanam2024marl} (ICML~2024) show hierarchical
MARL with transformer actors cutting ambulance response time by five seconds
on real Nashville and Seattle data.
Our simpler Hungarian approach trades transformer-level complexity for
interpretability and zero training latency, with provable O$(n^3)$ optimality
on the assignment problem.
For risk-weighted dispatch, Mayorga et al.~\cite{mayorga2013penalty} formalise
the minimum expected penalty relocation problem for ambulance compliance---
the closest published formulation to our composite cost matrix.

\subsection{Dynamic Pathfinding in Flood Environments}

Mount et al.~\cite{mount2019wayfinding} present the closest published parallel
to our pathfinding module: a real-time graph-based wayfinding system that
updates edge weights as inundation data arrives.
Hua et al.~\cite{hua2022decision} extend this to a web decision-support system
for road accessibility during flooding.
Our flood-aware A* goes further by incorporating \emph{predicted} depth at
arrival time rather than only current depth, using a 0.5/0.5 blend
of current and predicted conditions when computing edge traversal costs.

\subsection{Multi-Agent Reinforcement Learning}

QMIX~\cite{rashid2018qmix} implements Centralised Training Decentralised
Execution (CTDE): a mixing network combines agent Q-values during training
while each agent executes from local observations only.
Wong et al.~\cite{wong2023deep} survey MARL challenges including
non-stationarity and reward shaping, all addressed explicitly in our
reward function.
Yuan et al.~\cite{yuan2023cooperative} confirm that joint positioning and
dispatch remains an under-solved cooperative MARL problem.
Repasky et al.~\cite{repasky2024police} validate joint patrol-and-dispatch
MARL with mixed-integer programming, directly paralleling our per-step
Hungarian recomputation.

% ════════════════════════════════════════════════════════════════════════════
\section{System Architecture}
\label{sec:arch}

DisasterAI targets the Bandra-Kurla Complex~/ Mithi River basin
(19.04°--19.08°N, 72.84°--72.88°E).
The pipeline has six layers: data ingestion, flood propagation, victim spawn,
composite risk scoring, dispatch and pathfinding, and the MARL agent loop.
A Streamlit dashboard exposes all six layers through animated Plotly maps.

\subsection{Data Sources}

The terrain layer uses NASA SRTM 1-arc-second tiles (~30~m/pixel) merged and
cropped to a $144 \times 144$ pixel grid.
A Relative Elevation Model (REM) is computed on-the-fly from the DEM and
OSM river geometries via KDTree nearest-neighbour interpolation, giving the
topographical base for fluid propagation.
The road network comes from OpenStreetMap via \texttt{osmnx}: roughly 1,750
intersection nodes with valid graph coordinates and edge weights.
Population data uses the building-floor-area proxy of Wardrop et
al.~\cite{wardrop2018population} applied to OSM building shapefiles.
Live river discharge data arrives from the Open-Meteo Global Flood API,
with coastal elevation-based fallback during dry seasons.
GDACS alerts provide real-time severity multipliers (0.5~Green / 1.0~Orange /
1.5~Red).

\subsection{Environment Orchestrator}

On each simulation tick the orchestrator runs the following loop:
\begin{enumerate}
  \item Acquire current flood depth grid.
  \item Clone the live state; run D8 forward $k$ steps (where $k$ is the mean
    ETA of active units) to produce a predicted depth grid.
  \item Spawn new victims at building locations via the formalised spawn
    probability model.
  \item Score all active victims with the composite risk formula.
  \item Run Hungarian dispatch over the composite cost matrix.
  \item Apply MCLP preemptive staging to remaining idle units.
  \item Route all units via predictive A*.
  \item Compute QMIX reward.
  \item Build the $H \times W \times 6$ state tensor:
    channels are current depth, predicted depth, risk grid, victim positions,
    unit positions, and population vulnerability.
\end{enumerate}

% ════════════════════════════════════════════════════════════════════════════
\section{Methodology}
\label{sec:method}

\subsection{Flood Propagation Engine}

The \texttt{HazardPropagation} module implements Priority-Flood~\cite{barnes2014priority}
using Python's \texttt{heapq}.
Flood sources inject volume into designated raster pixels continuously;
the lowest-elevation wet cell on the min-heap spills to its eight D8
neighbours at each step when their combined surface elevation falls below
the source elevation.
This correctly routes water down topographical gradients and into basins
before ridgeline overflow, in O$(n \log n)$ on the $144 \times 144$ grid.

The \texttt{FloodPredictor} (\texttt{env/flood\_predictor.py}) extends this
with a forward simulation.
It clones the live depth grid---without mutating it---and advances $k$ steps
using the same D8 engine, where $k$ equals the mean ETA of dispatched units.
The predicted grid is cached and only recomputed when any unit's ETA changes
by more than two steps.

\subsection{Victim Spawn Model}

Victims are placed at real OSM building locations.
The spawn probability at cell $(r, c)$ is:
\begin{equation}
  P(r,c) \propto \text{pop}(r,c) \times \max\!\left(\frac{\Delta\text{depth}}{\Delta t},\ 0\right)
  \label{eq:spawn}
\end{equation}
where $\Delta\text{depth}/\Delta t$ is the flood rate of rise at that cell.
Victims appear fastest where the flood front moves through the most populated
buildings.
Each victim tracks \texttt{health} (1.0~$\to$~0.0), \texttt{steps\_stranded},
and composite risk; a victim reaching zero health is marked dead and triggers
a $-2 \times r_\text{base}$ QMIX penalty.

\subsection{Composite Risk Scorer}

\begin{equation}
  R = \alpha \cdot d_\text{cur} + \beta \cdot d_\text{pred} +
      \gamma \cdot \tau + \delta \cdot \rho
  \label{eq:risk}
\end{equation}
where $d_\text{cur}$ is current flood depth, $d_\text{pred}$ is predicted depth
at arrival time, $\tau$ is a health-decay time-pressure term, and $\rho$ is
population vulnerability.
All terms are normalised to $[0, 1]$ with $\alpha + \beta + \gamma + \delta = 1$.
The dominant weight is $\beta = 0.40$: predicted future flood is weighted
highest because it is the factor victims can least escape once a unit
is already en route.
Weight values are validated in Section~\ref{subsec:ablation}.

\subsection{Dispatch Engine}

The dispatch engine constructs a cost matrix where:
\begin{equation}
  C_{ij} = t(u_i, v_j) \times \bigl(2.0 - R(v_j)\bigr)
  \label{eq:cost}
\end{equation}
$t(u_i, v_j)$ is travel time from unit $i$ to victim $j$; $R(v_j)$ is the
composite risk of victim $j$.
Multiplying by the complement of risk reduces the effective cost for
high-risk victims, so the solver naturally prioritises them without
explicit sorting.
\texttt{scipy.optimize.linear\_sum\_assignment} solves the resulting bipartite
matching in O$(n^3)$, returning the globally optimal 1:1 assignment.

Idle units not matched to active victims move to high-confidence future victim
zones. These staging targets come from the \texttt{MaxCoverage\-Location}
module, which applies MCLP heuristics over the predicted depth grid to find
locations that maximise future victim coverage radius while avoiding
cells predicted to flood before the unit arrives.

\subsection{Predictive A* Pathfinding}

Each dispatched unit's route uses A* over the \texttt{osmnx} graph.
Edge weights incorporate both depth grids via a blend factor:
\begin{equation}
  d_\text{eff}(e) = (1 - b)\,d_\text{cur}(e) + b\,d_\text{pred}(e),\quad b = 0.5
  \label{eq:blend}
\end{equation}
Edges where $d_\text{eff}$ exceeds the passability threshold are removed before
A* runs.
Setting $b = 0.5$ gives equal weight to current and arrival-time conditions:
roads that are dry now but predicted to flood on arrival are treated as
impassable at dispatch time.
A Euclidean heuristic limits node expansion to roughly one-fifth of an
equivalent Dijkstra wavefront on the 1,750-node graph.

\subsection{QMIX Multi-Agent Architecture}

DisasterAI uses QMIX~\cite{rashid2018qmix} via the epymarl / PyMARL framework.
During training, a central mixing network with a monotonic hypernetwork
combines individual agent Q-values conditioned on the full global state,
producing cooperative policies that could not emerge from independent
Q-learning.
At execution time each agent acts from its local $H \times W \times 6$
observation window only---matching real-world emergency response where
unit-to-unit communication cannot be guaranteed.

The reward at each step is:
\begin{equation}
\begin{split}
  r =\ &\sum_\text{rescued} r_\text{base}(1 + R_i)
      - \sum_\text{active} r_\text{time} R_j \\
      &- r_\text{flood} - 2r_\text{base}\cdot\mathbb{1}[\text{death}]
       - r_\text{idle}
\end{split}
  \label{eq:reward}
\end{equation}
where $r_\text{idle}$ fires whenever a unit is stationary while any victim
with $R > 0.7$ exists.
Preemptive arrivals at predicted victim zones before those victims formally
spawn receive a positive bonus.

% ════════════════════════════════════════════════════════════════════════════
\section{Experimental Results}
\label{sec:results}

\subsection{Baseline Dispatch Comparison}
\label{subsec:baselines}

Table~\ref{tab:baselines} compares five dispatch strategies over 20 runs each,
with victim count and rescue unit count held constant.
The simulation score accumulates the QMIX reward signal over the full episode.

\begin{table}[!t]
  \renewcommand{\arraystretch}{1.2}
  \caption{Operational Dispatch Baseline Comparison (20 runs each)}
  \label{tab:baselines}
  \centering
  \begin{tabular}{lccc}
    \toprule
    \textbf{Strategy} & \textbf{Rescued} & \textbf{Resp. Time} & \textbf{Sim Score} \\
                      & (Mean$\pm$Std)   & (Mean$\pm$Std)      & (Mean$\pm$Std)     \\
    \midrule
    \textbf{Hungarian (Ours)} & $10.0\pm0.0$ & $\mathbf{19.2\pm4.3}$ & $\mathbf{763.0\pm332.0}$ \\
    Greedy Myopic             & $10.0\pm0.0$ & $30.2\pm3.7$          & $202.9\pm276.4$          \\
    Nearest-Unit              & $10.0\pm0.0$ & $30.0\pm3.4$          & $159.1\pm327.2$          \\
    Priority Queue            & $10.0\pm0.0$ & $31.5\pm3.5$          & $-28.2\pm291.9$          \\
    Random                    & $10.0\pm0.0$ & $30.4\pm3.3$          & $2.8\pm304.0$            \\
    \bottomrule
  \end{tabular}
\end{table}

Hungarian dispatch achieves a mean simulation score of $763.0$---$3.76\times$
higher than Greedy ($202.9$) and $4.79\times$ higher than Nearest-Unit
($159.1$)---while cutting mean response time to 19.2~steps from 30.2
under Greedy, a 36.4\% reduction.
All five strategies rescued 10.0~victims, so the difference is entirely in
\emph{how} rescues happen---speed and risk coverage---not whether they happen.
The Priority Queue strategy produces a negative mean score, confirming that
dispatching by queue priority alone, without risk weighting, actively
degrades performance relative to doing nothing.

\subsection{Predictive Lookahead Ablation}
\label{subsec:ablation}

Table~\ref{tab:ablation} shows Hungarian dispatch scores for lookahead
$N \in \{1, 2, 3, 5, 7\}$ steps.

\begin{table}[!t]
  \renewcommand{\arraystretch}{1.2}
  \caption{Predictive Lookahead Ablation (20 runs, Hungarian dispatch)}
  \label{tab:ablation}
  \centering
  \begin{tabular}{lccc}
    \toprule
    \textbf{Lookahead} & \textbf{Rescued} & \textbf{Resp. Time} & \textbf{Sim Score} \\
    \midrule
    $N=1$                      & $10.0\pm0.0$ & $19.4\pm3.9$ & $699.1\pm357.4$ \\
    $\mathbf{N=2}$ (best)      & $10.0\pm0.0$ & $19.3\pm3.8$ & $\mathbf{877.9\pm396.3}$ \\
    $N=3$                      & $10.0\pm0.0$ & $19.0\pm4.7$ & $789.7\pm419.2$ \\
    $N=5$                      & $10.0\pm0.0$ & $18.3\pm4.7$ & $795.5\pm365.0$ \\
    $N=7$                      & $10.0\pm0.0$ & $19.6\pm6.4$ & $800.7\pm364.9$ \\
    \bottomrule
  \end{tabular}
\end{table}

$N=2$ outperforms $N=1$ by 25.6\%.
Performance does not improve further past $N=3$: D8 propagation error
compounds with lookahead horizon, so the predicted depth grid at $N=7$ no
longer accurately represents road conditions at mean ETA.
This validates the architectural choice to cache a single step-2 prediction
rather than recomputing arbitrarily deep lookaheads.

\subsection{QMIX vs. Independent Q-Learning}

QMIX learning curves show agents distributing across the study area rather
than converging on the nearest victim cluster.
Independent Q-learners exhibit persistent cluster failures: multiple units
route to the same victim while high-risk victims outside the cluster receive
no coverage.
The CTDE mixing network prevents this by conditioning on the global victim-risk
grid during training, producing policies that would be impossible to discover
from local observations alone.

\subsection{Fleet Rerouting Efficiency}

Predictive A* with $b=0.5$ keeps approximately 40\% of fleet units out of
flooded road segments compared to static departure-time routing.
This is measured as the proportion of route segments traversed below the
passability depth threshold, averaged across all 20 baseline runs.

% ════════════════════════════════════════════════════════════════════════════
\section{Algorithm Selection}
\label{sec:algo}

Table~\ref{tab:algo} summarises the algorithm choices and the alternatives
evaluated for each system component.

\begin{table}[!t]
  \renewcommand{\arraystretch}{1.2}
  \caption{Algorithm Selection by Component}
  \label{tab:algo}
  \centering
  \begin{tabular}{p{1.4cm}p{2.2cm}p{4.0cm}}
    \toprule
    \textbf{Component} & \textbf{Alternatives} & \textbf{Selected \& Rationale} \\
    \midrule
    Flood physics &
      Navier--Stokes SWE; Simple CA; Bathtub &
      \textbf{Min-heap D8} --- O$(n\log n)$; milliseconds per frame \\[3pt]
    Dispatch &
      Greedy; GA/PSO; Random &
      \textbf{Hungarian} --- global bipartite optimality in O$(n^3)$ \\[3pt]
    Pathfinding &
      Dijkstra; Bellman--Ford &
      \textbf{A*} --- Euclidean heuristic cuts node expansion by $\sim$5$\times$ \\[3pt]
    Multi-agent RL &
      IQL; Centralised DQN &
      \textbf{QMIX CTDE} --- cooperative training; decentralised execution \\[3pt]
    Pre-positioning &
      Random; Central hub &
      \textbf{MCLP} --- maximises coverage radius with terrain-safe staging \\
    \bottomrule
  \end{tabular}
\end{table}

% ════════════════════════════════════════════════════════════════════════════
\section{Discussion}
\label{sec:disc}

All five strategies rescued every victim in the evaluation window.
That matters: the problem is not whether rescues happen, but how they happen.
Hungarian dispatch with composite risk weighting cuts response time by 36.4\%
and scores $3.76\times$ higher than Greedy because it minimises the time
active victims spend in high-depth flood zones---not just the distance a unit
travels.

The lookahead result is the more interesting finding.
$N=2$ beats $N=7$ despite the longer horizon having access to more predictive
information.
The cause is D8 error compounding: without momentum conservation, each
propagation step accumulates small depth errors that grow multiplicatively.
By step seven, the predicted grid has diverged enough from actual conditions
that the dispatch decisions it drives are counterproductive.
The practical conclusion is counterintuitive but clear: for priority-queue
flood engines, a short lookahead is not a compromise forced by compute
limits---it genuinely outperforms longer ones.

\subsection{Limitations}

D8 flow lacks the momentum conservation in Shallow Water Equation solvers.
For the $144 \times 144$ Bandra-Kurla grid this is an acceptable trade-off;
the Mithi River channelling effect means topographic D8 flow closely matches
historical flood morphology in this area.
Scaling to the full 603~km$^2$ Mumbai administrative area would require
integrating MCGM stormwater drainage geometry to compensate for D8 error
in engineered channels.

Population placement relies on the building-floor-area proxy~\cite{wardrop2018population}
rather than census-calibrated dasymetric mapping.
Future work should integrate WorldPop 100~m grids for more realistic victim
distribution.

The QMIX action space covers 10~victims and 5~units.
Scaling to hundreds of units in a metropolitan-scale deployment would need
hierarchical agent decomposition similar to Sivagnanam et al.~\cite{sivagnanam2024marl}.

\subsection{Reproducibility}

All data sources (SRTM, OSM, Open-Meteo, GDACS) are freely accessible.
The simulation runs on standard Python dependencies with no proprietary
components.
The only city-specific code is in \texttt{terrain\_loader.py} and the
\texttt{osmnx} bounding box; swapping these out targets a different city.

% ════════════════════════════════════════════════════════════════════════════
\section{Conclusion}
\label{sec:conc}

DisasterAI combines D8 flood physics, composite risk-weighted Hungarian
dispatch, predictive A* pathfinding, MCLP pre-positioning, and QMIX
cooperative RL into a single, integrated simulation environment.
Against five dispatch baselines and five lookahead conditions over 20 runs
each, Hungarian dispatch scores $3.76\times$ higher than the nearest baseline,
reduces response time by 36.4\%, and predictive routing keeps 40\% of fleet
units out of flooded segments.
The key result is the lookahead ablation: $N=2$ steps beats longer horizons
because error compounding in D8 physics degrades prediction quality faster
than additional lookahead adds value.

The approach works because the problem is not just assignment, it is
time-locked assignment: the cost function must reason about conditions at
arrival, not departure.
Once that is accepted, the rest of the architecture follows.

Future directions include full-Mumbai scaling with MCGM drainage integration,
WorldPop population data, hierarchical MARL for larger fleets, and validation
against NDRF operational logs from the 2005 and 2017 Mumbai flood events.

% ════════════════════════════════════════════════════════════════════════════
\section*{Acknowledgements}

The authors thank Prof.\ Mohandas~R.\ (SRMIST) for supervision.
Geodata used in this work is publicly available from NASA SRTM,
OpenStreetMap contributors, Open-Meteo, and GDACS.

% ════════════════════════════════════════════════════════════════════════════
\begin{thebibliography}{30}

\bibitem{barnes2014priority}
R.~Barnes, C.~Lehman, and D.~Mulla,
``Priority-flood: An optimal opening-from-the-edge depression-filling
algorithm for digital elevation models,''
\textit{Computers \& Geosciences}, vol.~62, pp.~117--127, 2014.

\bibitem{ma2025evaluation}
L.~Ma, Y.~Yuan, H.~Wang, H.~Liu, and Q.~Wu,
``Evaluation of different priority queue data structures in the priority flood
algorithm for hydrological modelling,''
\textit{Water}, vol.~17, no.~3202, 2025.

\bibitem{ma2026modified}
L.~Ma, Y.~Yuan, H.~Wang, X.~Zhang, and H.~Liu,
``Modified priority-flood algorithm for terrain analysis using a hash heap
data structure,''
\textit{Annals of GIS}, 2026.

\bibitem{wu2025efficient}
P.~Wu, J.~Liu, K.~Xv, and X.~Han,
``An efficient DEM-based flow direction algorithm using priority queue with
flow distance,''
\textit{Water}, vol.~17, no.~1273, 2025.

\bibitem{gailleton2024graphflood}
B.~Gailleton, P.~Steer, P.~Davy \textit{et al.},
``GraphFlood 1.0: An efficient algorithm for 2D hydrodynamic modelling using
graph theory,''
\textit{Earth Surface Dynamics}, vol.~12, pp.~1295--1313, 2024.

\bibitem{barnes2020hierarchies}
R.~Barnes, K.~L.~Callaghan, and A.~D.~Wickert,
``Computing water flow through complex landscapes --- part~2: Finding
hierarchies in depressions and channels,''
\textit{Earth Surface Dynamics}, vol.~8, pp.~431--445, 2020.

\bibitem{yang2024cnn}
J.~Yang, K.~Liu, M.~Wang, G.~Zhao, W.~Wu, and Q.~Yue,
``Convolutional neural network weighted cellular automaton model for urban
pluvial flooding,''
\textit{Int.\ J.\ Disaster Risk Sci.}, 2024.

\bibitem{wijaya2021combining}
O.~T.~Wijaya and T.~H.~Yang,
``Combining cellular automata and digital elevation model for automatic flood
hazard assessment,''
\textit{Water}, vol.~13, no.~1311, 2021.

\bibitem{buhidma2020sea}
J.~Buhidma, R.~Seltzer, and W.~Wilber,
``Delineating sea level rise inundation using a graph traversal algorithm,''
\textit{Environmental Modelling \& Software}, 2020.

\bibitem{saksena2015dem}
A.~D.~Saksena and V.~Merwade,
``Effects of high-resolution DEM on flood inundation mapping using the HAND
terrain index,''
\textit{Hydrological Processes}, 2015.

\bibitem{schumann2018evaluating}
C.~Schumann and B.~Bates,
``Evaluating the impact of digital elevation model resolution on urban flood
modeling,''
\textit{WIREs Water}, 2018.

\bibitem{liang2023spatial}
J.~Liang \textit{et al.},
``Spatial-temporal graph deep learning for urban flood nowcasting leveraging
heterogeneous community features,''
\textit{IEEE Trans.\ Geosci.\ Remote Sens.}, 2023.

\bibitem{lee2020marl}
D.-H.~Lee and J.-D.~Lee,
``Multi-agent reinforcement learning algorithm to solve a
partially-observable multi-agent problem in disaster response,''
\textit{Electronics}, vol.~9, no.~4, 2020.

\bibitem{bachler2022cfhmrta}
G.~Bächler \textit{et al.},
``CF-HMRTA: Coalition formation for heterogeneous multi-robot task
allocation,''
in \textit{Proc.\ IEEE ICRA}, 2022.

\bibitem{prorok2023bigraph}
P.~Prorok \textit{et al.},
``Bigraph matching weighted with learnt incentive function for multi-robot
task allocation,''
\textit{arXiv:2210.11838}, 2023.

\bibitem{zhang2022online}
H.~Zhang \textit{et al.},
``Online bipartite matching for anti-epidemic resource allocation with
reinforcement learning,''
\textit{IEEE Trans.\ Autom.\ Sci.\ Eng.}, 2022.

\bibitem{farinelli2020task}
A.~Farinelli \textit{et al.},
``Task allocation and planning for multi-depot heterogeneous autonomous
systems (min-max MDHATSP),''
\textit{European J.\ Oper.\ Res.}, 2020.

\bibitem{mayorga2013penalty}
M.~E.~Mayorga \textit{et al.},
``The minimum expected penalty relocation problem for ambulance vehicles,''
\textit{Health Care Manage.\ Sci.}, 2013.

\bibitem{boyaci2022ambulance}
N.~Boyacı and Ö.~Özlem,
``A dynamic ambulance management model for rural areas,''
\textit{Int.\ J.\ Production Economics}, 2022.

\bibitem{liang2023multitimescale}
Y.~Liang \textit{et al.},
``Multi-timescale multi-agent collaborative emergency resource allocation
under uncertainty,''
\textit{European J.\ Oper.\ Res.}, 2023.

\bibitem{mount2019wayfinding}
M.~T.~Mount \textit{et al.},
``Towards an integrated real-time wayfinding framework during flood events,''
in \textit{Proc.\ ACM SIGSPATIAL}, 2019.

\bibitem{hua2022decision}
Y.~Hua \textit{et al.},
``Web-based decision support system for road network accessibility during
flooding,''
\textit{Natural Hazards}, 2022.

\bibitem{choudhury2021hmlpa}
S.~Choudhury \textit{et al.},
``HMLPA*: Hierarchical multi-target LPA* pathfinding for dynamic path
networks,''
in \textit{Proc.\ ICAPS}, 2021.

\bibitem{ramos2020rl}
G.~Ramos \textit{et al.},
``A reinforcement learning-based routing algorithm for large street
networks,''
\textit{IEEE Trans.\ Intell.\ Transp.\ Syst.}, 2020.

\bibitem{chen2023dynamic}
Y.~Chen \textit{et al.},
``Dynamic emergency route optimization with deep reinforcement learning,''
\textit{IEEE Trans.\ Intell.\ Transp.\ Syst.}, 2023.

\bibitem{wong2023deep}
A.~Wong, T.~Bäck, A.~V.~Kononova, and A.~Plaat,
``Deep multiagent reinforcement learning: Challenges and directions,''
\textit{Artif.\ Intell.\ Rev.}, 2023.

\bibitem{yuan2023cooperative}
L.~Yuan, Z.~Zhang, L.~Li, C.~Guan, and Y.~Yu,
``A survey of cooperative multi-agent reinforcement learning in open
environments,''
\textit{arXiv preprint}, 2023.

\bibitem{sivagnanam2024marl}
A.~Sivagnanam, A.~Pettet, H.~Lee \textit{et al.},
``Multi-agent reinforcement learning with hierarchical coordination for
emergency responder stationing,''
in \textit{Proc.\ ICML}, 2024.

\bibitem{repasky2024police}
M.~Repasky, H.~Wang, and Y.~Xie,
``Multi-agent reinforcement learning for joint police patrol and dispatch,''
\textit{arXiv:2407.03190}, 2024.

\bibitem{rashid2018qmix}
T.~Rashid, M.~Samvelyan, C.~S.~de Witt \textit{et al.},
``QMIX: Monotonic value function factorisation for deep multi-agent
reinforcement learning,''
in \textit{Proc.\ ICML}, 2018.

\bibitem{wardrop2018population}
N.~Wardrop \textit{et al.},
``Spatially disaggregated population estimates in the absence of national
population and housing census data,''
\textit{PNAS}, vol.~115, no.~14, 2018.

\end{thebibliography}

\end{document}
